\section*{Finite source-resolved error calculus for projected multilinear DAGs}
\label{sec:finite-source-resolved-error-calculus}

\paragraph{Scope.}
Let $G=(V,E)$ be a finite acyclic typed graph.  Every internal vertex $v$
has a finite-dimensional multilinear law $\mu_v$, and every type carries an
orthogonal projector $P_v$ (the projector may be the same for several
vertices).  A leaf has a projected baseline $R_v$ and a finite sum of labelled
local source vectors.  A closure residual at an internal vertex is treated as
one additional labelled degree-one source with amplitude one.  All source
labels are finite, and no infinite series or universal sharpness assertion is
part of this theorem.

\paragraph{The finite polynomial theorem.}
For a node $v$, write
\[
  D_v(t)=\sum_{\alpha\ne0} A_{v,\alpha}t^\alpha
\]
for its raw error relative to the recursively projected baseline, where
multi-indices record source multiplicities.  If the ordered children have
baselines $R_1,\ldots,R_a$ and errors $D_1,\ldots,D_a$, multilinearity gives
\[
\begin{split}
 D_v(t)={}&
 \sum_{\varnothing\ne S\subseteq\{1,\ldots,a\}}
 \mu_v(z^S_1,\ldots,z^S_a)
 +(I-P_v)\mu_v(R_1,\ldots,R_a)\\
 &+L_v(t),
\end{split}
\tag{1}
\]
where $z_i^S=D_i$ for $i\in S$, $z_i^S=R_i$ otherwise, and $L_v$ is the
explicit local-source polynomial.  At a leaf, $D_v=L_v$.  Since each child
polynomial is finite and each $\mu_v$ is multilinear, (1) is a finite
polynomial identity.  Coefficients are obtained by Cartesian convolution of
the child supports, merging indices by addition; consequently a source used
twice has multiplicity two rather than being silently identified with a
degree-one source.

\paragraph{Proof.}
Expand
\[
 \mu_v(R_1+D_1,\ldots,R_a+D_a)
\]
by multilinearity.  The all-baseline term is
$\mu_v(R_1,\ldots,R_a)$.  The recursively projected baseline is
$P_v\mu_v(R_1,\ldots,R_a)$, so subtracting it leaves the normal residual
$(I-P_v)\mu_v(R_1,\ldots,R_a)$.  Every other term chooses a nonempty subset
of erroneous slots, giving the first sum in (1); adding $L_v$ gives the final
term.  Induction over a topological ordering proves the assertion at every
node.  The same induction proves that a cached topological evaluator and a
recursive evaluator that unrolls shared dependencies have identical
coefficients.  At the projected root,
\[
 P_r(I-P_r)\mu_r(R_1,\ldots,R_a)=0,
\]
so the root-local normal source is removed exactly by the output projection.
\qed

\paragraph{First-order source recombination.}
After retaining only degree-one terms, let $A_{v,s}$ be the coefficient
operator of source $s$.  If $H_{v,u}$ is the linearized edge operator and
$L_{v,s}$ the local source operator, then
\[
 A_{v,s}=L_{v,s}+\sum_{u\to v}H_{v,u}A_{u,s}.
\tag{2}
\]
Thus all paths carrying the same label are added before taking a norm.  With
source vectors $\epsilon_s$,
\[
 \left\|\sum_s A_{r,s}\epsilon_s\right\|
 \le B_{\rm source}:=\sum_s\|A_{r,s}\|\,\|\epsilon_s\|
 \le B_{\rm path},
\tag{3}
\]
where $B_{\rm path}$ takes the norm separately on every path.  Equation (2)
is a topological dynamic program; its graph traversal is $O(|V|+|E|)$ apart
from matrix multiplication costs.

\paragraph{Finite truncation.}
For any order $p$, split $D_r=D_r^{(\le p)}+R_r^{(>p)}$.  For amplitudes
$t_s$ the exact finite polynomial satisfies
\[
 \|R_r^{(>p)}(t)\|
 \le \sum_{|\alpha|>p}\|A_{r,\alpha}\|
                 \prod_s|t_s|^{\alpha_s}.
\tag{4}
\]
No omitted term is discarded; (4) is the triangle bound over the finite
omitted support.

\paragraph{Signed compositional expressions.}
Let $F=\sum_j c_jT_j$ and let
$P_j(t)=\sum_{\alpha\ne0}A_{j,\alpha}t^\alpha$ be the finite source
polynomial of $T_j$.  Define
\[
 A_{F,\alpha}=\sum_j c_jA_{j,\alpha},\qquad
 B_{\rm signed}=\sum_\alpha\|A_{F,\alpha}\|
                              \prod_s|t_s|^{\alpha_s}.
\]
Then
\[
 \left\|\sum_jc_jP_j(t)\right\|
 \le B_{\rm signed}
 \le B_{\rm treewise}:=
 \sum_j|c_j|\sum_\alpha\|A_{j,\alpha}\|
                              \prod_s|t_s|^{\alpha_s}.
\tag{5}
\]
The first inequality is the triangle inequality after coefficient
aggregation; the second is coefficientwise subadditivity.  Repeated-source
and mixed-source terms remain distinct multi-indices throughout.  Applying
(5) to a truncation and its omitted support gives the corresponding signed
remainder certificate.

\paragraph{Current boundary.}
The theorem is exact for finite declared DAGs and finite source supports.  It
does not provide a scalable compression theorem for implicitly generated
graphs, an infinite-series tail envelope, a universal signed-forest sharp
constant, or theorem-level novelty.  Those boundaries remain separately
registered rather than being inferred from the finite calculus.
